\section{Лекция 5. Числа Каталана и линейные рекуррентные последовательности}
\noeqref{eq:psp_prefix, eq:catalan_diff, eq:catalan_closed, eq:catalan_conv}
\noeqref{eq:catalan_ratio, eq:catalan_series, eq:catalan_part_sum}
\noeqref{eq:catalan_asymp, eq:rec_seq_def, eq:lrs_def, eq:char_poly}
\noeqref{eq:lrs_gen_sol, eq:lrs_ord2_diff, eq:lrs_ord2_mult, eq:binet_formula}
\noeqref{eq:phidias_limit}

\subsection{Числа Каталана и правильные скобочные последовательности}

\subsubsection{Мотивирующая задача об очереди в кассу}

\begin{problem}[\textbf{Задача о кассе и сдаче}]
    В кассу за билетами выстроилась очередь из $2n$ покупателей. Известно, 
    что билет стоит $500$ рублей, причём ровно у $n$ покупателей есть 
    только купюра достоинством $500$ рублей, а у остальных $n$ покупателей 
    --- только купюра достоинством $1000$ рублей. В начальный момент в 
    кассе нет денег (размена нет). Сколькими способами $2n$ покупателей 
    могут расположиться в очереди (если различать только тип купюры 
    покупателя), чтобы кассир всегда мог выдать сдачу каждому покупателю с 
    купюрой в $1000$ рублей без задержек?
\end{problem}

\begin{figure}[H]
    \centering
    \begin{tikzpicture}[scale=0.9]
        % Касса справа
        \draw[thick, fill=gray!10] (6.5, 0) rectangle (7.8, 1.4);
        \draw[thick, fill=gray!25] (6.3, 1.4) -- (7.15, 2.0) -- (8.0, 1.4) -- cycle;
        \draw[thick, fill=white] (6.85, 0.4) rectangle (7.45, 1.0);
        \node[below, align=center] at (7.15, -0.1) 
            {\textbf{Касса}\\($500$ рублей)};

        % Покупатели в очереди
        % Покупатель 1 (ближний к кассе)
        \draw[thick] (5.0, 1.1) circle (0.2);
        \draw[thick] (5.0, 0.9) -- (5.0, 0.35);
        \draw[thick] (4.75, 0.7) -- (5.25, 0.7);
        \draw[thick] (5.0, 0.35) -- (4.8, 0);
        \draw[thick] (5.0, 0.35) -- (5.2, 0);
        \node[below] at (5.0, -0.05) {$500$};

        % Покупатель 2
        \draw[thick] (3.6, 1.1) circle (0.2);
        \draw[thick] (3.6, 0.9) -- (3.6, 0.35);
        \draw[thick] (3.35, 0.7) -- (3.85, 0.7);
        \draw[thick] (3.6, 0.35) -- (3.4, 0);
        \draw[thick] (3.6, 0.35) -- (3.8, 0);
        \node[below] at (3.6, -0.05) {$500$};

        % Многоточие
        \node at (2.2, 0.6) {$\dots$};

        % Покупатель 2n (последний в очереди)
        \draw[thick] (0.8, 1.1) circle (0.2);
        \draw[thick] (0.8, 0.9) -- (0.8, 0.35);
        \draw[thick] (0.55, 0.7) -- (1.05, 0.7);
        \draw[thick] (0.8, 0.35) -- (0.6, 0);
        \draw[thick] (0.8, 0.35) -- (1.0, 0);
        \node[below] at (0.8, -0.05) {$1000$};

        % Фигурная скобка снизу
        \draw[thick] (0.4, -0.65) -- (0.4, -0.8) -- (5.4, -0.8) -- (5.4, -0.65);
        \node[below] at (2.9, -0.8) {Всего $2n$ человек ($n$ по $500$ и $n$ по $1000$)};
    \end{tikzpicture}
    \caption{Модель очереди из $2n$ покупателей у билетной кассы.}
\end{figure}

\begin{remark}[Формализация задачи через скобочные последовательности]
    Сопоставим каждому покупателю с купюрой $500$ рублей открывающую 
    скобку «$($» (он увеличивает запас пятисотрублёвых купюр в кассе на $1$), 
    а каждому покупателю с купюрой $1000$ рублей --- закрывающую скобку 
    «$)$» (ему требуется выдать одну пятисотрублёвую купюру в качестве 
    сдачи, что уменьшает запас кассы на $1$). Тогда кассир сможет обслужить 
    всех покупателей тогда и только тогда, когда в любой момент времени 
    число вошедших покупателей с $500$ рублями не меньше числа обслуженных 
    покупателей с $1000$ рублями.
\end{remark}

\begin{definition}
    Слово длины $2n$ в двухбуквенном алфавите $\{ (, ) \}$ называется 
    \textdef{правильной скобочной последовательностью (ПСП)}, если 
    одновременно выполняются два условия:
    \begin{enumerate}
        \item Общее количество открывающих скобок «$($» равно общему 
        количеству закрывающих скобок «$)$» (и составляет по $n$ штук);
        \item В любом начальном отрезке (\textdef{префиксе}) слова 
        количество открывающих скобок не меньше количества закрывающих:
        \begin{equation} \label{eq:psp_prefix}
            \forall k \in \{1, 2, \dots, 2n\} : 
            \#\text{«}(\text{» в префиксе длины } k \ge 
            \#\text{«})\text{» в префиксе длины } k. \tag{1}
        \end{equation}
    \end{enumerate}
\end{definition}

\begin{definition}
    \textdef{Числом Каталана порядка $n$} (обозначается $C_n$) называется 
    количество всех правильных скобочных последовательностей длины $2n$. 
    По определению при $n = 0$ полагают $C_0 = 1$ (единственная ПСП длины 
    $0$ --- пустое слово $\eps$).
\end{definition}

\subsubsection{Геометрический вывод формулы чисел Каталана}

Поставим в соответствие каждой скобочной последовательности длины $2n$ 
ломаную линию (\textdef{траекторию}) на координатной плоскости $Oxy$, 
выходящую из начала координат $(0, 0)$:
\begin{itemize}
    \item открывающей скобке «$($» сопоставим шаг вправо-вверх на вектор 
    $(+1, +1)$;
    \item закрывающей скобке «$)$» сопоставим шаг вправо-вниз на вектор 
    $(+1, -1)$.
\end{itemize}

Если в последовательности ровно $n$ открывающих и $n$ закрывающих скобок, 
то через $2n$ шагов траектория обязательно закончится в точке $(2n, 0)$.
\begin{itemize}
    \item \textdef{«Хорошая» траектория} (\textdef{путь Дика}) соответствует 
    ПСП: она целиком лежит не ниже оси абсцисс ($y \ge 0$), хотя может 
    касаться оси $Ox$.
    \item \textdef{«Плохая» траектория} соответствует последовательности, 
    нарушающей условие \eqref{eq:psp_prefix}: хотя бы в одной точке она 
    опускается строго ниже оси $Ox$ ($y < 0$), а значит, хотя бы один раз 
    пересекает или касается горизонтальной прямой $y = -1$.
\end{itemize}

\begin{figure}[H]
    \centering
    \begin{tikzpicture}[scale=0.68]
        % Оси координат
        \draw[->, thick] (-0.5, 0) -- (11.2, 0) node[below] {$x$};
        \draw[->, thick] (0, -2.2) -- (0, 4.2) node[left] {$y$};
        \node[below left] at (0, 0) {$0$};

        % Деления по осям
        \foreach \x in {1,...,10} {
            \draw (\x, 0.1) -- (\x, -0.1) node[below, scale=0.8] {$\x$};
        }
        \foreach \y in {-1, 1, 2, 3, 4} {
            \draw (0.1, \y) -- (-0.1, \y) node[left, scale=0.8] {$\y$};
        }

        % Сетка
        \draw[gray!25, very thin] (0, -2) grid (10, 4);

        % Хорошая траектория: ( ( ) ) ( ( ) ( ) )
        \draw[blue!80!black, very thick] 
            (0,0) -- (1,1) -- (2,2) -- (3,1) -- (4,0) -- (5,1) 
            -- (6,2) -- (7,1) -- (8,2) -- (9,1) -- (10,0);
        \foreach \x/\y in {0/0, 1/1, 2/2, 3/1, 4/0, 5/1, 6/2, 7/1, 8/2, 9/1, 10/0} {
            \fill[blue!80!black] (\x, \y) circle (2.2pt);
        }

        % Плохой участок траектории: ( ( ) ) ) ( ( ( ) )
        \draw[red!80!black, very thick] (4,0) -- (5,-1) -- (6,0) -- (7,1);
        \fill[red!80!black] (5, -1) circle (2.5pt);
        \fill[red!80!black] (6, 0) circle (2.2pt);

        % Подписи скобочных последовательностей
        \node[blue!80!black, font=\bfseries] at (5, 3.6) {ПСП («хорошая» траектория):};
        \node[blue!80!black] at (5, 2.9) 
            {$(\quad (\quad )\quad )\quad (\quad (\quad )\quad (\quad )\quad )$};

        \node[red!80!black] at (5, -1.6) 
            {$(\quad (\quad )\quad )\quad \mathbf{)}\quad (\quad (\quad (\quad )\quad )$};
        \node[red!80!black, font=\bfseries] at (5, -2.2) 
            {Не ПСП («плохая» траектория: опускается до $y = -1$)};
    \end{tikzpicture}
    \caption{Геометрическая интерпретация скобочных последовательностей длины 
    $2n = 10$: «хорошая» траектория (синяя) не опускается ниже оси $Ox$, 
    а «плохая» (красная) достигает уровня $y = -1$.}
\end{figure}

\begin{theorem}[\textbf{Явная формула для чисел Каталана}]
    Для любого целого $n \ge 0$ число Каталана $C_n$ вычисляется по формуле:
    \begin{equation} \label{eq:catalan_closed}
        C_n = C_{2n}^n - C_{2n}^{n+1} = \frac{1}{n + 1} C_{2n}^n = 
        \frac{(2n)!}{n!(n + 1)!} \tag{2}
    \end{equation}
\end{theorem}

\begin{proof}[Доказательство (принцип отражения Андре)]~\\
    \begin{quote}
        Обозначим через $N_{\text{хор}} = C_n$ число «хороших» траекторий 
        из точки $(0, 0)$ в точку $(2n, 0)$, через $N_{\text{общ}}$ --- 
        общее число всех траекторий из $(0, 0)$ в $(2n, 0)$, а через 
        $N_{\text{плох}}$ --- число «плохих» траекторий. Очевидно, что:
        \[
            C_n = N_{\text{хор}} = N_{\text{общ}} - N_{\text{плох}}.
        \]

        \textbf{1. Подсчёт общего числа траекторий $N_{\text{общ}}$:}\\
        Любая траектория из $(0, 0)$ в $(2n, 0)$ состоит ровно из $n$ шагов 
        вверх «$($» и $n$ шагов вниз «$)$». Выбрать $n$ позиций для 
        открывающих скобок из $2n$ мест можно $C_{2n}^n$ способами, поэтому:
        \[
            N_{\text{общ}} = C_{2n}^n.
        \]

        \textbf{2. Подсчёт числа «плохих» траекторий $N_{\text{плох}}$ 
        (принцип отражения):}\\
        Поскольку каждая «плохая» траектория опускается ниже оси $Ox$, а 
        ордината за один шаг меняется на $\pm 1$, на любой «плохой» 
        траектории существует \emph{первая} (самая левая) точка $M(k, -1)$ 
        касания прямой $y = -1$.

        Возьмём участок «плохой» траектории от начальной точки $(0, 0)$ до 
        первой точки $M(k, -1)$ на прямой $y = -1$ и \textbf{отразим его 
        симметрично относительно горизонтальной прямой $y = -1$}, а 
        оставшуюся часть траектории от точки $M(k, -1)$ до $(2n, 0)$ 
        оставим без изменений (см. рисунок ниже).

        \begin{figure}[H]
            \centering
            \begin{tikzpicture}[scale=0.68]
                % Оси координат
                \draw[->, thick] (-0.5, 0) -- (11.2, 0) node[below] {$x$};
                \draw[->, thick] (0, -4.5) -- (0, 3.0) node[left] {$y$};
                \node[below left] at (0, 0) {$0$};

                % Деления
                \foreach \x in {1,...,10} {
                    \draw (\x, 0.1) -- (\x, -0.1) node[above, scale=0.75] {$\x$};
                }
                \foreach \y in {-4, -3, -2, -1, 1, 2} {
                    \draw (0.1, \y) -- (-0.1, \y) node[left, scale=0.75] {$\y$};
                }

                % Сетка
                \draw[gray!25, very thin] (0, -4) grid (10, 2);

                % Ось отражения y = -1
                \draw[red, dashed, thick] (0, -1) -- (10.5, -1) 
                    node[right, scale=0.85] {ось отражения $y = -1$};

                % Исходный начальный участок (пунктиром)
                \draw[gray!80!black, dashed, thick] 
                    (0,0) -- (1,1) -- (2,2) -- (3,1) -- (4,0) -- (5,-1);

                % Отражённый начальный участок из (0, -2) до (5, -1)
                \draw[red!80!black, very thick] 
                    (0,-2) -- (1,-3) -- (2,-4) -- (3,-3) -- (4,-2) -- (5,-1);
                \foreach \x/\y in {0/-2, 1/-3, 2/-4, 3/-3, 4/-2} {
                    \fill[red!80!black] (\x, \y) circle (2.2pt);
                }

                % Оставшаяся часть траектории от (5, -1) до (10, 0)
                \draw[blue!80!black, very thick] 
                    (5,-1) -- (6,0) -- (7,1) -- (8,2) -- (9,1) -- (10,0);
                \fill[red!80!black] (5, -1) circle (3pt) 
                    node[below right, scale=0.85] {$M(5, -1)$};
            \end{tikzpicture}
            \caption{Принцип отражения: симметричное отражение начального 
            участка «плохой» траектории относительно прямой $y = -1$ 
            переводит точку старта $(0, 0)$ в точку $(0, -2)$.}
        \end{figure}

        При симметричном отражении относительно прямой $y = -1$ начало 
        координат $(0, 0)$ переходит в точку $(0, -2)$, поэтому всякая 
        «плохая» траектория однозначно превращается в некоторую 
        траекторию из точки $(0, -2)$ в точку $(2n, 0)$.

        Обратно, любая траектория, идущая из точки $(0, -2)$ в точку 
        $(2n, 0)$, начинается ниже прямой $y = -1$, а заканчивается выше 
        неё, поэтому по дискретной непрерывности она обязана пересечь 
        прямую $y = -1$ в некоторой первой точке $M$. Отражая участок от 
        $(0, -2)$ до $M$ обратно относительно $y = -1$, мы единственным 
        образом восстанавливаем исходную «плохую» траекторию.

        Таким образом, установлено биективное соответствие между всеми 
        «плохими» траекториями из $(0, 0)$ в $(2n, 0)$ и \emph{всеми} 
        вообще траекториями из $(0, -2)$ в $(2n, 0)$. Чтобы пройти за $2n$ 
        шагов из уровня $y = -2$ на уровень $y = 0$ (подняться в сумме на 
        $\Delta y = +2$), траектория должна содержать ровно $n + 1$ шагов 
        вверх и $n - 1$ шагов вниз:
        \[
            \begin{cases}
                u + d = 2n, \\
                u - d = 2
            \end{cases}
            \implies u = n + 1, \quad d = n - 1.
        \]
        Следовательно, число «плохих» траекторий равно:
        \[
            N_{\text{плох}} = C_{2n}^{n+1} = C_{2n}^{n-1}.
        \]

        \textbf{3. Вычисление разности:}\\
        Вычитая $N_{\text{плох}}$ из $N_{\text{общ}}$ и преобразуя 
        факториалы, получаем формулу \eqref{eq:catalan_closed}:
        \begin{equation} \label{eq:catalan_diff}
            \begin{aligned}
                C_n &= C_{2n}^n - C_{2n}^{n+1} = 
                \frac{(2n)!}{n! \, n!} - \frac{(2n)!}{(n + 1)!(n - 1)!} = 
                C_{2n}^n - \frac{n}{n + 1} C_{2n}^n \\
                &= \frac{1}{n + 1} C_{2n}^n = \frac{(2n)!}{n!(n + 1)!}.
            \end{aligned} \tag{3}
        \end{equation}
        Теорема доказана. Первые числа Каталана при $n = 0, 1, 2, 3, 4, 5$ 
        равны соответственно: $1, 1, 2, 5, 14, 42$.
    \end{quote}
\end{proof}

\subsection{Свойства чисел Каталана}

\subsubsection{Свойство 1. Рекуррентное соотношение Сегнера (свёртка)}

\begin{theorem}[\textbf{Свёртка чисел Каталана}]
    Для любого натурального $n \ge 1$ числа Каталана удовлетворяют 
    нелинейному рекуррентному соотношению:
    \begin{equation} \label{eq:catalan_conv}
        C_n = \sum_{k=0}^{n-1} C_k \cdot C_{n-k-1}, \qquad C_0 = 1. \tag{4}
    \end{equation}
\end{theorem}

\begin{proof}[Доказательство]~\\
    \begin{quote}
        Рассмотрим произвольную непустую правильную скобочную 
        последовательность $w$ длины $2n$ ($n \ge 1$). Её первый символ 
        обязательно является открывающей скобкой «$($». Поскольку $w$ --- 
        ПСП, у этой первой открывающей скобки существует единственная 
        парная ей закрывающая скобка «$)$» (геометрически это первая точка 
        $(2k + 2, 0)$, в которой траектория Дика впервые возвращается на 
        ось $Ox$).

        Следовательно, любая ПСП $w$ длины $2n$ \emph{единственным образом} 
        представима в каноническом виде:
        \[
            w = (w_1) w_2,
        \]
        где $w_1$ и $w_2$ --- две независимые правильные скобочные 
        последовательности (возможно, пустые). Пусть подпоследовательность 
        $w_1$ содержит $k$ пар скобок, то есть имеет длину $2k$, где индекс 
        $k$ может принимать любое значение из множества $\{0, 1, \dots, n - 1\}$. 
        Тогда на долю подпоследовательности $w_2$ остаётся длина:
        \[
            2n - 2 - 2k = 2(n - k - 1),
        \]
        то есть $n - k - 1$ пар скобок. При фиксированном $k$ выбрать ПСП 
        $w_1$ можно $C_k$ способами, а ПСП $w_2$ --- независимо $C_{n-k-1}$ 
        способами. По правилу умножения и правилу сложения (суммируя по 
        всем $k = 0, 1, \dots, n - 1$) получаем формулу \eqref{eq:catalan_conv}.
    \end{quote}
\end{proof}

\subsubsection{Свойство 2. Связь соседних чисел Каталана}

\begin{proposition}
    Для любого натурального $n \ge 1$ выполняется соотношение:
    \begin{equation} \label{eq:catalan_ratio}
        C_n = \frac{2(2n - 1)}{n + 1} C_{n-1} \tag{5}
    \end{equation}
\end{proposition}

\begin{proof}[Доказательство]~\\
    \begin{quote}
        Воспользуемся явной факториальной формулой \eqref{eq:catalan_closed} 
        и распишем отношение $C_n$ к $C_{n-1}$:
        \begin{align*}
            C_n &= \frac{(2n)!}{n!(n + 1)!} = 
            \frac{2n \cdot (2n - 1) \cdot (2n - 2)!}
            {n \cdot (n - 1)! \cdot (n + 1) \cdot n!} \\
            &= \frac{2 \cancel{n} (2n - 1)}{\cancel{n}(n + 1)} \cdot 
            \frac{(2(n - 1))!}{(n - 1)! \, n!} = 
            \frac{2(2n - 1)}{n + 1} C_{n-1}.
        \end{align*}
        Соотношение \eqref{eq:catalan_ratio} позволяет быстро вычислять 
        числа Каталана последовательно одно за другим за $O(n)$ 
        арифметических операций.
    \end{quote}
\end{proof}

\subsubsection{Свойство 3. Сумма ряда с числами Каталана}

\begin{proposition}
    Справедливо равенство для суммы бесконечного числового ряда:
    \begin{equation} \label{eq:catalan_series}
        \sum_{k=0}^{\infty} \frac{C_k}{4^k} = 2 \tag{6}
    \end{equation}
\end{proposition}

\begin{proof}[Доказательство (Способ 1: через частичные суммы по индукции)]~\\
    \begin{quote}
        Докажем предварительно по индукции, что для любого $n \ge 0$ 
        конечная частичная сумма ряда выражается тождеством:
        \begin{equation} \label{eq:catalan_part_sum}
            S_n = \sum_{k=0}^{n} \frac{C_k}{4^k} = 
            2 - \frac{2n + 1}{4^n} C_n = 2 - \frac{n + 1}{2 \cdot 4^n} C_{n+1}. 
            \tag{7}
        \end{equation}
        \begin{itemize}
            \item \emph{База индукции ($n = 0$):} 
            $S_0 = \frac{C_0}{4^0} = 1$, и по правой части 
            $2 - \frac{1}{1} \cdot 1 = 1$. База верна.
            \item \emph{Индукционный шаг ($n - 1 \implies n$):} 
            Пусть формула \eqref{eq:catalan_part_sum} верна для $n - 1$. 
            Прибавим к $S_{n-1}$ слагаемое $\frac{C_n}{4^n}$ и выразим 
            $C_{n-1}$ через $C_n$ из свойства \eqref{eq:catalan_ratio} 
            ($\frac{2n - 1}{4^{n-1}} C_{n-1} = \frac{2(n + 1)}{4^n} C_n$):
            \[
                S_n = S_{n-1} + \frac{C_n}{4^n} = 
                2 - \frac{2n - 1}{4^{n-1}} C_{n-1} + \frac{C_n}{4^n} = 
                2 - \frac{2n + 2}{4^n} C_n + \frac{C_n}{4^n} = 
                2 - \frac{2n + 1}{4^n} C_n.
            \]
        \end{itemize}
        Используя асимптотику чисел Каталана $C_n \sim \frac{4^n}{\sqrt{\pi n^3}}$ 
        (доказываемую ниже в свойстве 4), перейдём к пределу при $n \to \infty$:
        \[
            \lim_{n \to \infty} \frac{2n + 1}{4^n} C_n = 
            \lim_{n \to \infty} \frac{2n + 1}{\cancel{4^n}} \cdot 
            \frac{\cancel{4^n}}{\sqrt{\pi} n^{3/2}} = 
            \lim_{n \to \infty} \frac{2}{\sqrt{\pi n}} = 0.
        \]
        Следовательно, $\sum_{k=0}^\infty \frac{C_k}{4^k} = 
        \lim_{n \to \infty} S_n = 2 - 0 = 2$.
    \end{quote}
\end{proof}

\begin{proof}[Доказательство (Способ 2: через производящую функцию)]~\\
    \begin{quote}
        Рассмотрим производящую функцию последовательности чисел Каталана 
        $c(x) = \sum_{n=0}^\infty C_n x^n$. Умножая свёртку 
        \eqref{eq:catalan_conv} на $x^n$ и суммируя по $n \ge 1$, получаем 
        квадратное уравнение:
        \[
            c(x) - 1 = x \cdot \bigl(c(x)\bigr)^2 \iff 
            x c^2(x) - c(x) + 1 = 0.
        \]
        Выбирая ветвь корня, удовлетворяющую условию $c(0) = C_0 = 1$, 
        находим явный вид производящей функции:
        \[
            c(x) = \frac{1 - \sqrt{1 - 4x}}{2x}.
        \]
        Подставляя $x = \frac{1}{4}$ (по первой теореме Абеля о 
        непрерывности суммы степенного ряда вплоть до границы круга 
        сходимости), получаем:
        \[
            \sum_{k=0}^{\infty} \frac{C_k}{4^k} = c\left(\frac{1}{4}\right) = 
            \frac{1 - \sqrt{1 - 4 \cdot \frac{1}{4}}}{2 \cdot \frac{1}{4}} = 
            \frac{1 - 0}{1/2} = 2.
        \]
    \end{quote}
\end{proof}

\subsubsection{Свойство 4. Асимптотика чисел Каталана}

\begin{theorem}[\textbf{Асимптотическая формула для чисел Каталана}]
    При $n \to \infty$ для чисел Каталана справедливо асимптотическое 
    соотношение:
    \begin{equation} \label{eq:catalan_asymp}
        C_n \sim \frac{4^n}{\sqrt{\pi n^3}} \tag{8}
    \end{equation}
\end{theorem}

\begin{proof}[Доказательство]~\\
    \begin{quote}
        Запишем число Каталана в виде $C_n = \frac{1}{n + 1} \frac{(2n)!}{(n!)^2}$ 
        и применим к факториалам $(2n)!$ и $n!$ формулу Стирлинга 
        $m! \sim \sqrt{2 \pi m} \left(\frac{m}{e}\right)^m$:
        \begin{align*}
            (2n)! &\sim \sqrt{2 \pi (2n)} \left(\frac{2n}{e}\right)^{2n} = 
            2\sqrt{\pi n} \cdot \frac{4^n n^{2n}}{e^{2n}}, \\
            (n!)^2 &\sim \left(\sqrt{2 \pi n} \left(\frac{n}{e}\right)^n\right)^2 = 
            2 \pi n \cdot \frac{n^{2n}}{e^{2n}}.
        \end{align*}
        Подставим эти асимптотические выражения в формулу для $C_n$:
        \[
            C_n = \frac{1}{n + 1} \frac{(2n)!}{(n!)^2} \sim 
            \frac{1}{n + 1} \cdot 
            \frac{\cancel{2}\sqrt{\pi n} \cdot 4^n \cancel{n^{2n} / e^{2n}}}
            {\cancel{2} \pi n \cdot \cancel{n^{2n} / e^{2n}}} = 
            \frac{1}{n + 1} \cdot \frac{4^n}{\sqrt{\pi n}}.
        \]
        Так как при $n \to \infty$ выполняется предел $\lim_{n \to \infty} 
        \frac{n + 1}{n} = \lim_{n \to \infty} \left(1 + \frac{1}{n}\right) = 1$, 
        множитель $n + 1$ в знаменателе эквивалентен $n$ ($n + 1 \sim n$):
        \[
            C_n \sim \frac{4^n}{n \sqrt{\pi n}} = \frac{4^n}{\sqrt{\pi n^3}}.
        \]
        Формула \eqref{eq:catalan_asymp} доказана.
    \end{quote}
\end{proof}

\subsection{Линейные рекуррентные последовательности (ЛРП)}

\subsubsection{Основные определения и постановка задачи}

\begin{definition}
    Пусть $\{x_i\}_{i=1}^\infty$ --- числовая последовательность. Говорят, 
    что $\{x_i\}_{i=1}^\infty$ является \textdef{рекуррентной (возвратной) 
    последовательностью порядка $k$}, если существует функция $F$ от $k$ 
    переменных, такая что начиная с номера $n > k$ каждый член 
    последовательности выражается через $k$ предыдущих:
    \begin{equation} \label{eq:rec_seq_def}
        x_n = F(x_{n-1}, x_{n-2}, \dots, x_{n-k}) \qquad \forall n > k. \tag{9}
    \end{equation}
\end{definition}

\begin{definition}
    Возвратная последовательность порядка $k$ называется \textdef{линейной 
    рекуррентной последовательностью (ЛРП)} (однородной, с постоянными 
    коэффициентами), если функция $F$ линейна:
    \begin{equation} \label{eq:lrs_def}
        x_n = \alpha_1 x_{n-1} + \alpha_2 x_{n-2} + \dots + \alpha_k x_{n-k}, 
        \qquad n > k, \tag{10}
    \end{equation}
    где $\alpha_1, \alpha_2, \dots, \alpha_k \in \C$ --- фиксированные 
    коэффициенты, причём $\alpha_k \neq 0$.
\end{definition}

\begin{remark}[Задача о поиске явной формулы $x_n = f(n)$]
    Чтобы однозначно выделить конкретную последовательность, удовлетворяющую 
    рекуррентному уравнению $k$-го порядка \eqref{eq:lrs_def}, необходимо 
    задать $k$ начальных значений (\textdef{начальных условий}):
    \[
        x_1 = a_1, \quad x_2 = a_2, \quad \dots, \quad x_k = a_k.
    \]
    Основная задача теории ЛРП состоит в том, чтобы по рекуррентному 
    соотношению \eqref{eq:lrs_def} и начальным условиям найти замкнутую 
    аналитическую формулу общего члена $x_n = f(n)$, зависящую только от 
    номера $n$.
\end{remark}

\subsubsection{Характеристический многочлен и теорема об общем решении}

Будем искать частные решения уравнения \eqref{eq:lrs_def} в виде 
геометрической прогрессии $x_n = z^n$ ($z \neq 0$). Подставляя $x_n = z^n$ 
в \eqref{eq:lrs_def} и сокращая на $z^{n-k}$, получаем алгебраическое 
уравнение $k$-й степени:
\[
    z^n = \alpha_1 z^{n-1} + \alpha_2 z^{n-2} + \dots + \alpha_k z^{n-k} \iff 
    z^k - \alpha_1 z^{k-1} - \dots - \alpha_{k-1} z - \alpha_k = 0.
\]

\begin{definition}
    Многочлен степени $k$ от комплексной переменной $z \in \C$:
    \begin{equation} \label{eq:char_poly}
        P(z) = z^k - \alpha_1 z^{k-1} - \dots - \alpha_{k-1} z - \alpha_k \tag{11}
    \end{equation}
    называется \textdef{характеристическим многочленом} линейной 
    рекуррентной последовательности \eqref{eq:lrs_def}, а уравнение 
    $P(z) = 0$ --- её \textdef{характеристическим уравнением}.
\end{definition}

По основной теореме алгебры характеристический многочлен $P(z)$ степени $k$ 
имеет в поле комплексных чисел $\C$ ровно $k$ корней с учётом их кратностей. 
Пусть $\lambda_1, \lambda_2, \dots, \lambda_s$ --- все попарно различные 
корни многочлена $P(z)$ с кратностями $r_1, r_2, \dots, r_s$ соответственно, 
так что:
\[
    P(z) = (z - \lambda_1)^{r_1} (z - \lambda_2)^{r_2} \dots (z - \lambda_s)^{r_s}, 
    \qquad \sum_{i=1}^{s} r_i = k.
\]

\begin{theorem}[\textbf{Об общем решении однородной ЛРП}]
    Общее решение линейного рекуррентного уравнения \eqref{eq:lrs_def} 
    порядка $k$ задаётся явно формулой:
    \begin{equation} \label{eq:lrs_gen_sol}
        x_n = \sum_{i=1}^{s} \left( c_{1,i} + c_{2,i} \, n + c_{3,i} \, n^2 + 
        \dots + c_{r_i, i} \, n^{r_i - 1} \right) \lambda_i^n, \tag{12}
    \end{equation}
    где $c_{j,i} \in \C$ ($i = \overline{1, s}$, $j = \overline{1, r_i}$) --- 
    произвольные постоянные (всего $\sum_{i=1}^s r_i = k$ констант), которые 
    однозначно определяются из системы $k$ линейных алгебраических 
    уравнений, порождённой начальными условиями:
    \[
        \begin{cases}
            x_1 = a_1, \\
            x_2 = a_2, \\
            \dots \\
            x_k = a_k.
        \end{cases}
    \]
\end{theorem}

\subsubsection{Расчётные формулы для ЛРП 2-го порядка}

Рассмотрим наиболее важный на практике случай линейной рекуррентной 
последовательности второго порядка ($k = 2$):
\[
    x_n = \alpha_1 x_{n-1} + \alpha_2 x_{n-2}, \qquad x_1 = a_1, \quad x_2 = a_2.
\]
Характеристическое уравнение является квадратным:
\[
    P(z) = z^2 - \alpha_1 z - \alpha_2 = 0, \qquad 
    D = \alpha_1^2 + 4\alpha_2.
\]
В зависимости от дискриминанта $D$ возникают два случая:
\begin{enumerate}
    \item \textbf{Случай простых корней ($D \neq 0$, то есть 
    $\lambda_1 \neq \lambda_2$, кратности $r_1 = r_2 = 1$):} 
    Общая формула \eqref{eq:lrs_gen_sol} принимает вид линейной комбинации 
    двух геометрических прогрессий:
    \begin{equation} \label{eq:lrs_ord2_diff}
        x_n = c_1 \lambda_1^n + c_2 \lambda_2^n, \tag{13}
    \end{equation}
    а константы $c_1, c_2$ находятся из системы двух линейных уравнений:
    \[
        \begin{cases}
            c_1 \lambda_1 + c_2 \lambda_2 = a_1, \\
            c_1 \lambda_1^2 + c_2 \lambda_2^2 = a_2.
        \end{cases}
    \]
    \item \textbf{Случай кратного корня ($D = 0$, один корень 
    $\lambda_1 = \frac{\alpha_1}{2}$ кратности $r_1 = 2$):} 
    Согласно теореме \eqref{eq:lrs_gen_sol} общее решение имеет вид:
    \begin{equation} \label{eq:lrs_ord2_mult}
        x_n = (c_1 + c_2 n) \lambda_1^n, \tag{14}
    \end{equation}
    где константы $c_1, c_2$ находятся из системы:
    \[
        \begin{cases}
            (c_1 + c_2) \lambda_1 = a_1, \\
            (c_1 + 2 c_2) \lambda_1^2 = a_2.
        \end{cases}
    \]
\end{enumerate}

\subsubsection{Последовательность Фибоначчи и вывод формулы Бине}

\begin{example}[\textbf{Вывод формулы Бине для чисел Фибоначчи}] ~
    Найдём явную формулу общего члена последовательности Фибоначчи 
    $\{f_n\}_{n=1}^\infty$, заданной рекуррентным соотношением второго 
    порядка и начальными условиями:
    \[
        f_n = f_{n-1} + f_{n-2} \quad (n \ge 3), \qquad f_1 = 1, \quad f_2 = 1.
    \]

    \begin{proof}[Решение]~\\
        \begin{quote}
            \textbf{1. Корни характеристического многочлена:}\\
            Перенесём члены в одну часть ($f_n - f_{n-1} - f_{n-2} = 0$) и 
            выпишем характеристический многочлен \eqref{eq:char_poly}:
            \[
                P(z) = z^2 - z - 1 = 0.
            \]
            Дискриминант равен $D = (-1)^2 - 4 \cdot 1 \cdot (-1) = 5 > 0$. 
            Характеристическое уравнение имеет два различных вещественных 
            корня (кратности $r_1 = r_2 = 1$):
            \[
                \lambda_1 = \frac{1 + \sqrt{5}}{2}, \qquad
                \lambda_2 = \frac{1 - \sqrt{5}}{2}.
            \]
            Следовательно, по формуле \eqref{eq:lrs_ord2_diff} общий член 
            последовательности Фибоначчи ищется в виде:
            \[
                f_n = c_1 \lambda_1^n + c_2 \lambda_2^n.
            \]

            \textbf{2. Нахождение констант $c_1$ и $c_2$ из начальных 
            условий:}\\
            Подставим начальные условия $f_1 = 1$ и $f_2 = 1$:
            \[
                \begin{cases}
                    f_1 = c_1 \lambda_1 + c_2 \lambda_2 = 1, \\
                    f_2 = c_1 \lambda_1^2 + c_2 \lambda_2^2 = 1
                \end{cases}
                \iff
                \begin{cases}
                    c_1 \left(\dfrac{1 + \sqrt{5}}{2}\right) + 
                    c_2 \left(\dfrac{1 - \sqrt{5}}{2}\right) = 1, \\[3mm]
                    c_1 \left(\dfrac{1 + \sqrt{5}}{2}\right)^2 + 
                    c_2 \left(\dfrac{1 - \sqrt{5}}{2}\right)^2 = 1.
                \end{cases}
            \]
            Чтобы решить эту систему максимально просто, заметим, что 
            поскольку $\lambda_1, \lambda_2$ являются корнями уравнения 
            $z^2 = z + 1$ (либо продолжая последовательность назад до 
            нулевого члена $f_0 = f_2 - f_1 = 0$), вычитание первого 
            уравнения из второго сразу даёт условие для $n = 0$:
            \[
                c_1(\lambda_1^2 - \lambda_1) + c_2(\lambda_2^2 - \lambda_2) = 
                1 - 1 \implies c_1 + c_2 = 0 \implies c_2 = -c_1.
            \]
            Подставим $c_2 = -c_1$ в первое уравнение системы:
            \[
                c_1 (\lambda_1 - \lambda_2) = 1 \implies 
                c_1 \left( \frac{1 + \sqrt{5}}{2} - \frac{1 - \sqrt{5}}{2} \right) = 1 
                \implies c_1 \sqrt{5} = 1.
            \]
            Отсюда находим коэффициенты:
            \[
                c_1 = \frac{1}{\sqrt{5}}, \qquad c_2 = -\frac{1}{\sqrt{5}}.
            \]

            \textbf{3. Формула Бине:}\\
            Подставляя найденные значения $c_1$ и $c_2$, получаем знаменитую 
            \textdef{формулу Бине}:
            \begin{equation} \label{eq:binet_formula}
                f_n = \frac{1}{\sqrt{5}} \left(\frac{1 + \sqrt{5}}{2}\right)^n - 
                \frac{1}{\sqrt{5}} \left(\frac{1 - \sqrt{5}}{2}\right)^n = 
                \frac{\left(\dfrac{1 + \sqrt{5}}{2}\right)^n - 
                \left(\dfrac{1 - \sqrt{5}}{2}\right)^n}{\sqrt{5}}. \tag{15}
            \end{equation}
        \end{quote}
    \end{proof}
\end{example}

\begin{remark}[\textbf{Число Фидия (золотое сечение) и асимптотика чисел Фибоначчи}]
    Положительный корень характеристического уравнения чисел Фибоначчи:
    \[
        \Phi = \lambda_1 = \frac{1 + \sqrt{5}}{2} \approx 1,6180339887\dots
    \]
    называется \textdef{числом Фидия} (или \textdef{золотым сечением}). 
    Поскольку второй корень по модулю меньше единицы 
    ($\abs{\lambda_2} = \frac{\sqrt{5} - 1}{2} \approx 0,618 < 1$), при 
    $n \to \infty$ величина $\lambda_2^n \to 0$ стремится к нулю 
    экспоненциально быстро. Благодаря этому отношение соседних чисел 
    Фибоначчи в пределе стремится к числу Фидия:
    \begin{equation} \label{eq:phidias_limit}
        \lim_{n \to \infty} \frac{f_{n+1}}{f_n} = 
        \lim_{n \to \infty} \frac{\lambda_1^{n+1} - \lambda_2^{n+1}}
        {\lambda_1^n - \lambda_2^n} = 
        \lambda_1 \lim_{n \to \infty} 
        \frac{1 - (\lambda_2 / \lambda_1)^{n+1}}{1 - (\lambda_2 / \lambda_1)^n} = 
        \lambda_1 = \frac{1 + \sqrt{5}}{2} = \Phi. \tag{16}
    \end{equation}
\end{remark}